\documentclass[3p]{elsarticle}

\biboptions{super,sort&compress}

% Standard package includes for math, tables, and images
\usepackage{float}
\usepackage{enumitem}
\usepackage{subcaption}
\usepackage{amsmath}
\usepackage{booktabs}
\usepackage{graphicx}
\usepackage{url}
\usepackage{stfloats}
\usepackage{tabularx}
\usepackage{setspace}
\usepackage{xcolor}

% Font and layout improvements
\usepackage[T1]{fontenc}
\usepackage{lmodern} % scalable T1 fonts (needed for microtype on TeX installs without cm-super)
\usepackage[utf8]{inputenc}
\usepackage{microtype}
\usepackage{inconsolata}
\usepackage{hyperref}

% ---------------------------------------------------------------------------
% Placeholders for results that do not exist yet. Every number produced by the
% fills them, in paper/RESULTS_TODO.md. Do NOT replace a key with a number by
% hand: copy it from the named output file.
% Set \showTBDkeystrue to print the keys in table cells as well.
% ---------------------------------------------------------------------------
\newif\ifshowTBDkeys
\showTBDkeysfalse
\newcommand{\TBD}[1]{\textcolor{red}{[\textbf{TBD}: \nolinkurl{#1}]}}
\newcommand{\TBDtext}[2]{\textcolor{red}{[\textbf{TBD} \nolinkurl{#1}: \emph{#2}]}}
\newcommand{\TBDcell}[1]{\ifshowTBDkeys\textcolor{red}{\scriptsize\texttt{\detokenize{#1}}}\else\textcolor{red}{TBD}\fi}
\newcommand{\FIGTBD}[2]{\fbox{\parbox[c][4.5cm][c]{0.9\linewidth}{\centering\textcolor{red}{\textbf{Figure placeholder}}\\[2pt]\texttt{\detokenize{#1}}\\[2pt]\small #2}}}

\journal{Digital Discovery (manuscript prepared for submission)}

\emergencystretch=3em
\begin{document}

\begin{frontmatter}

    \title{Publication-aware prediction of battery material properties from literature-extracted data}

    \author[1]{Sunraj Patel}
    \author[1]{Shikhar Misra\corref{cor1}}

    \cortext[cor1]{Corresponding author.}
    \ead{shikharm@iitk.ac.in}

    \address[1]{Department of Materials Science and Engineering, Indian Institute of Technology Kanpur, India 208016}


    \begin{abstract}
        Text-mined databases provide hundreds of thousands of battery-material property records and are an attractive training source for machine-learning (ML) models, but their records are clustered by publication and evaluation protocols rarely account for it. We trained CatBoost--XGBoost ensembles for five properties (capacity, voltage, $\log_{10}$ conductivity, Coulombic efficiency and energy) on 285,243 records of the auto-generated Huang--Cole database and evaluated identical models under record-level, publication-grouped (DOI), temporal and composition-grouped splits. Between 30\% and 84\% of target variance lies between publications, yet grouping by DOI lowered $R^2$ only modestly (capacity \TBD{m/capacity/random/t2/r2} to \TBD{m/capacity/doi/t2/r2}). Three analyses explain the difference. Composition descriptors carry little publication identity: a model given only journal and year predicted from them has $R^2 \approx 0$, and composition predicts the journal only when publications are shared between training and test. A further leakage route survives DOI grouping: \TBD{rep/doi/share-range} of DOI-grouped test records repeat a (material, value) pair reported by another publication in training, and these are predicted far better ($R^2$ \TBD{rep/capacity/doi/r2_rep} for capacity) than values without such a precedent (\TBD{rep/capacity/doi/r2_new}). Holding out whole compositions lowers performance further (capacity \TBD{m/capacity/material/t2/r2}). Leakage therefore needs both clustered targets and features that can recognise the cluster, and a DOI-grouped split removes only one of several routes. We release an auditable pipeline with an explicit list of excluded columns, physical-range filters, a treatment of the database's warning flags, and adapted re-runs of the two previous regression pipelines on the same database.
    \end{abstract}

    \begin{keyword}
        Battery materials; Text mining; ChemDataExtractor; Materials informatics; Publication-grouped validation; Data leakage; Uncertainty.
    \end{keyword}

\end{frontmatter}

% ---------------- INTRODUCTION ----------------
\section{Introduction}
\label{sec:intro}

The performance of a rechargeable battery is bounded by the properties of its electrode and electrolyte materials. Specific capacity and the average electrode potential together set the specific energy of a cell; the ionic conductivity of the electrolyte limits rate capability; and the Coulombic efficiency measures how much charge is lost to side reactions in each cycle. New electrode and electrolyte chemistries are reported at a rate of thousands of papers per year, and each paper typically reports a handful of these quantities for one or a few materials, under conditions such as current density and cycle number that the authors choose.

Text-mining tools such as ChemDataExtractor \cite{swain2016chemdataextractor,mavracic2021chemdataextractor2} make it possible to turn this literature into a structured database. Huang and Cole applied a battery-specific version of ChemDataExtractor to 229,061 papers and obtained 292,313 records of five properties: capacity, voltage, conductivity, Coulombic efficiency and energy \cite{huang2020database,huang2020figshare}. A database of this size is an obvious candidate for training property-prediction models, in the same way that manually curated experimental databases have been used for photovoltaic and other device data. Earlier ML work on battery materials was limited by data volume: Kauwe et al.\ trained regressors on a manually digitised set of a few hundred cathode records \cite{kauwe2019battery,ghadbeigi2015performance}, and Huang and Cole explicitly motivated their database as a way of overcoming that limitation \cite{huang2020database}.

Literature-derived records, however, are clustered by publication. A single paper may report capacities at several current densities, for several cycles, for the parent material and its doped or carbon-coated variants, all prepared and measured in the same laboratory and reported in the same style. Huang and Cole flag 117,403 records (about 40\% of the database) as members of such within-paper data series \cite{huang2020database}. If records are split into training and test sets at random, most test records have siblings from the same paper in the training set, and a model can score well by recognising the paper rather than by learning a transferable relationship between a material and its properties. Grouped, cluster-based and time-split validation are established remedies for this kind of dependence \cite{roberts2017crossvalidation,meredig2018cluster,durdy2022locoCV,witman2025matfold,sheridan2013timesplit,kapoor2023leakage}. For a literature-derived photovoltaic database, a model's $R^2$ fell substantially when it was tested on records registered later rather than on a random subset \cite{fukasawa2024effectiveness}, and Jablonka recently showed that composition descriptors can predict the authors, journal and year of a record, so that models may succeed for bibliographic rather than chemical reasons \cite{jablonka2026clever}.

Here we ask how much this structure affects the evaluation of models trained on an auto-generated battery database, which is text-mined rather than curated, in long format (one property value per record), large, and almost free of experimental-condition fields. The answer is not the one a simple leakage argument predicts. The battery records are strongly clustered by publication, but the record-level and DOI-grouped estimates differ only modestly, and we trace this to two mechanisms: the features available here identify publications poorly, and many reported values recur across publications, so that a DOI-grouped split still allows a model to recall them. The workflow is summarised in Fig.~\ref{fig:workflow}. Our contributions are:

\begin{enumerate}[label=(\roman*)]
    \item \textbf{Multi-target, multi-protocol evaluation.} For each of five properties we compare record-level random, DOI-grouped, temporal and composition-grouped splits with the same model, features and partition fractions, and we quantify how much of each target's variance lies between publications.

    \item \textbf{Leakage channels beyond publication identity.} We show that composition features carry little recoverable publication identity, measure how much of the test signal comes from values that recur across publications, and separate theoretical from measured capacities.

    \item \textbf{Auditable feature construction and controls.} We list every excluded column with its reason, document physical-range filters and the treatment of the database's warning flags (one of which encodes publication identity and another the target), and report noise, paper-level-feature, identity and within-paper-series controls.

    \item \textbf{Comparison on a common basis.} We adapt the public pipeline of a previous study on the same database to our partitions and compare it with our ensemble on identical test records.
\end{enumerate}

\begin{figure}[htbp]
    \centering
    \FIGTBD{figures/battery_workflow.png}{Workflow: battery.csv $\rightarrow$ cleaning and physics filters $\rightarrow$ per-target feature tiers $\rightarrow$ random / DOI-grouped / temporal splits $\rightarrow$ CatBoost--XGBoost ensemble $\rightarrow$ metrics, grouped SHAP, ensemble-spread uncertainty.}
    \caption{\textbf{Workflow for publication-aware prediction of battery material properties.} Records from the Huang--Cole database are cleaned, filtered to physically plausible ranges and split by property. Composition, elemental-property, material-name and measurement-context descriptors form the main feature set; paper-level title keywords, the raw material name and Gaussian-noise columns are evaluated as separate controls. Each target is evaluated under record-level random, DOI-grouped and forward-in-time splits.}
    \label{fig:workflow}
\end{figure}

Most data-driven battery studies fall into two groups: cell-level models trained on cycling data from a controlled experimental campaign, and materials-level models trained on property values collected from the literature.

Severson et al.\ cycled 124 commercial LiFePO$_4$/graphite cells under fast-charging protocols and showed that features computed from the first 100 cycles predicted cycle life with a test error of 9.1\%; a classifier using the first five cycles separated short- from long-lived cells with 4.9\% error \cite{severson2019cyclelife}. Their cells were divided into a training set (41 cells), a primary test set (43 cells) and a secondary test set (40 cells) that was generated after model development, so the secondary test was a genuinely prospective evaluation at the level of whole cells.
Celik et al.\ used early-cycle data from 121 commercial LiFePO$_4$/graphite cells gathered from the literature and reported a cycle-life accuracy of 99.81\% with an extremely randomised trees model, and they attributed a considerable reduction in error to data management and physical understanding of the features \cite{celik2022cyclelife}. They used a 75/25 training/testing division; whether the division was made by cell or over cycles is not stated.
In both studies the unit of observation is a cell, all cells share one chemistry and manufacturer, and the question of publication-level dependence does not arise in the same form.

At the materials level, Kauwe, Rhone and Sparks predicted the specific discharge capacity at cycle 25 of cathode materials from the manually digitised dataset of Ghadbeigi et al.\ \cite{kauwe2019battery,ghadbeigi2015performance}. After removing columns with excessive missing entries and descriptors correlated with the target, their training data had 215 entries and 49 composition-derived and battery-metadata descriptors. They compared lasso, kernel ridge regression, random forests and support vector regression; the three non-linear models reached test $R^2$ of 0.48--0.51 and an MAE of about 16~mAh~g$^{-1}$, using an 80/20 training/test split, 10-fold cross-validation for model validation, and hyperparameters chosen by grid search against the test error \cite{kauwe2019battery}. No grouping of records by source is described.
The authors attributed the modest performance to missing descriptors, non-standard cell architectures and testing protocols, and unreported electrode composition \cite{kauwe2019battery}. The same reporting heterogeneity is present, at much larger scale, in a text-mined database.

ChemDataExtractor is a rule- and grammar-based toolkit for extracting chemical entities and property relations from scientific text \cite{swain2016chemdataextractor}; version 2.0 added nested, ontology-driven property models \cite{mavracic2021chemdataextractor2}. Huang and Cole adapted version~1.5 to battery papers and released the database used here \cite{huang2020database,huang2020figshare}. The authors estimated a precision of 80.0\% and a recall of 59.1\% from a manually evaluated sample, with property-level precision between 75.5\% (energy) and 83.3\% (capacity) \cite{huang2020database}. They later used fine-tuned BatteryBERT language models to classify papers and electrode roles and to enhance the database \cite{huang2022batterybert}. AtomWork-Battery takes a different route: records on solid electrolytes and cathode materials are extracted and curated by trained editors from papers selected by NLP screening, and are linked to crystal-structure and computed electronic-structure databases \cite{xu2024atomworkbattery}. More recently, large language models have been used for structured extraction of materials data \cite{dagdelen2024structured}; the questions of evaluation design raised below apply equally to databases built that way.

Few studies have trained models on the Huang--Cole records themselves. Zhuang and Barnard used its formulas to study structure-free composition encodings, in unsupervised clustering and classification of the resulting clusters \cite{zhuang2023mendeleev,zhuang2024classification,zhuang2025applications}, and trained XGBoost regressors on 254,749 cleaned records with a random 85/15 split, reporting test $R^2$ of 0.47 for capacity, 0.49 for voltage, 0.50 for energy and below 0.15 for conductivity and Coulombic efficiency \cite{zhuang2024property}. Jablonka used its capacity records as one of five case studies showing that composition descriptors also predict bibliographic metadata (authors, journal and year), and argued that group and time splits should be routine falsification tests \cite{jablonka2026clever}; that study evaluated its battery models by random cross-validation. Huang and Cole themselves cite Kauwe et al.\ as prior ML work, but that study used the earlier manually digitised dataset \cite{huang2020database,kauwe2019battery}. We screened all works citing the database (Section~\ref{sec:previous}); applications that do not cite it may exist.

Random cross-validation estimates performance on new records drawn from the same distribution as the training records. When observations are grouped, by site, patient, batch or publication, it estimates performance on new members of groups already seen, which is usually not the question of interest \cite{roberts2017crossvalidation,kapoor2023leakage}. In materials informatics, Meredig et al.\ introduced leave-one-cluster-out cross-validation (LOCO-CV) to measure extrapolation to chemically distinct materials \cite{meredig2018cluster}; Durdy et al.\ proposed kernelised LOCO-CV and random-projection baselines \cite{durdy2022locoCV}; Witman and Schindler standardised a hierarchy of increasingly strict chemical and structural hold-outs in the MatFold package \cite{witman2025matfold}. Related work examined extrapolation to extraordinary or high-$T_c$ materials \cite{stanev2018superconductor,kauwe2020extraordinary} and forward cross-validation ordered by property value \cite{xiong2020forward}. Time-split validation, in which models are trained on older and tested on newer compounds, estimates prospective performance \cite{sheridan2013timesplit}. For a literature-derived photovoltaic database, Fukasawa et al.\ reported that $R^2$ fell from 0.72 under a random split to 0.45 on records registered later \cite{fukasawa2024effectiveness}.

These approaches group by chemistry, by property value or by time. Grouping by source publication, which is the natural unit of dependence in a text-mined database, is less common. Of the battery ML studies discussed above, none reports a publication-grouped evaluation: the cycling studies use data from one campaign, Kauwe et al.\ used a small curated set, and the studies on the Huang--Cole database used random splits or random cross-validation, although Jablonka explicitly recommends group and time splits \cite{jablonka2026clever}. This is a gap in the studies we examined, not evidence that publication-grouped evaluation has not been done elsewhere.

% ---------------- METHODS ----------------
\section{Methods}

\subsection{Data Preparation}

\subsubsection{Dataset Description}
\label{sec:data}

We used the un-merged file \texttt{battery.csv} from the Huang--Cole release \cite{huang2020figshare}, which contains one row per extracted record: 292,313 records from 42,873 DOIs (Table~\ref{tab:dataset}). The release also contains a merged file in which identical (compound, property, value) triples from different papers are collapsed into one row with a record count. Merging removes the link between a record and its publication, which is required for grouped evaluation, so the merged file was not used.

The data are in long format. Each record carries one value of one property (\texttt{Property}, \texttt{Value}, normalised \texttt{Unit}, and the text as written in \texttt{Raw\_value} and \texttt{Raw\_unit}), the compound name as it appears in the text (\texttt{Name}), its parsed composition (\texttt{Extracted\_name}, a list of element--stoichiometry dictionaries), the property specifier recognised by the parser (\texttt{Specifier}), and publication metadata (\texttt{DOI}, \texttt{Title}, \texttt{Journal}, \texttt{Date}). Three further fields are sparse: \texttt{Info} stores the current density and cycle number found next to a capacity value; \texttt{Type} records whether the material was identified as a cathode, anode or electrolyte and is populated for about 3\% of records; and \texttt{Correctness} holds the manual evaluation labels for the validation sample of about 500 records \cite{huang2020database}. The \texttt{Tag} field marks 6,718 energy records as \emph{Calculated}: these were derived by Huang and Cole from capacity and voltage records of the same material rather than extracted from text. The \texttt{Warning} field is described in Section~\ref{sec:warnings}.

Two properties of this structure matter for what follows. First, there is no record that describes a cell or an experiment as a whole: the capacity and voltage of the same material in the same paper are separate records, and they cannot in general be joined because the database does not identify which values were measured together. Each property is therefore modelled separately. Second, the predictors available for a record are limited to what can be read from the compound name, its composition, the specifier and a few conditions. Electrode formulation, loading, electrolyte, counter electrode, voltage window and temperature, which strongly influence measured capacity and Coulombic efficiency, are not recorded.

\begin{table}[htbp]
    \centering\small
    \caption{\textbf{Records in the Huang--Cole database and in the modelling sets.} Raw counts are from \texttt{battery.csv}; publications are unique DOIs after normalisation. Cleaned counts are after the filters in Section~\ref{sec:cleaning}.}
    \label{tab:dataset}
    \begin{tabular}{@{}lrrrr@{}}
        \toprule
        Property & Raw records & Raw publications & Modelling records & Modelling publications \\
        \midrule
        Capacity               & 144,359 & \TBDcell{pop/capacity/raw_publications} & \TBDcell{pop/capacity/records} & \TBDcell{pop/capacity/publications} \\
        Voltage                & 114,240 & \TBDcell{pop/voltage/raw_publications} & \TBDcell{pop/voltage/records} & \TBDcell{pop/voltage/publications} \\
        Energy                 &  15,543 & \TBDcell{pop/energy/raw_publications} & \TBDcell{pop/energy/records} & \TBDcell{pop/energy/publications} \\
        Coulombic efficiency   &  11,003 & \TBDcell{pop/coulombic_efficiency/raw_publications} & \TBDcell{pop/coulombic_efficiency/records} & \TBDcell{pop/coulombic_efficiency/publications} \\
        Conductivity           &   7,168 & \TBDcell{pop/conductivity/raw_publications} & \TBDcell{pop/conductivity/records} & \TBDcell{pop/conductivity/publications} \\
        \midrule
        Total                  & 292,313 & 42,873 & \TBDcell{pop/total/records} & \TBDcell{pop/total/publications} \\
        \midrule
        \multicolumn{5}{@{}l}{\textbf{Flags in the raw file}} \\
        \multicolumn{3}{@{}l}{\texttt{Tag} = Calculated (energy only)} & \multicolumn{2}{r}{6,718} \\
        \multicolumn{3}{@{}l}{\texttt{Warning} = S / L / LS / R / RL} & \multicolumn{2}{r}{99,671 / 33,858 / 17,732 / 7,999 / 3,338} \\
        \multicolumn{3}{@{}l}{\texttt{Type} populated} & \multicolumn{2}{r}{$\approx$3\%} \\
        \multicolumn{3}{@{}l}{\texttt{Correctness} populated (manual sample)} & \multicolumn{2}{r}{$\approx$500} \\
        \bottomrule
    \end{tabular}
\end{table}

\subsubsection{Targets}
\label{sec:targets}

One regression model was trained per property (Table~\ref{tab:targets}). Conductivity was modelled as $\log_{10}(\sigma / \mathrm{S\,cm^{-1}})$ because reported values span many orders of magnitude and the physically meaningful error is multiplicative; its MAE and RMSE are therefore in decades. All other targets were modelled in their normalised units. A small fraction of capacity records carry the unit Ah~kg$^{-1}$, which is numerically identical to mAh~g$^{-1}$, and were kept.

The five targets are not equally well defined. ``Capacity'' pools theoretical, initial, reversible, charge and discharge capacities at different current densities and cycle numbers; ``voltage'' pools average, plateau, open-circuit and cut-off voltages; ``conductivity'' pools ionic and electronic conductivities of electrolytes and electrodes. The specifier and the parsed conditions are given to the model as features (Section~\ref{sec:features}), but a large part of the target variance is due to what was measured rather than to the material, and this sets a ceiling on achievable accuracy.

\subsection{Literature Corpus and Automated Extraction}

\subsubsection{Huang--Cole corpus}

The corpus is the output of the battery-specific ChemDataExtractor pipeline built by Huang and Cole \cite{huang2020database,huang2020figshare}, which was applied to 229,061 papers. The database stores each extracted value as a separate record, together with the compound name, its parsed elemental composition, the property specifier (for example ``discharge capacity'' or ``theoretical capacity''), additional conditions for capacity (current and cycle), the DOI, title, journal and publication date. This study does not perform any extraction of its own, and no large-language-model extraction was used: all information about the literature comes from the released records, and the released Correctness labels and Warning flags are the only quality information available. Later extraction approaches \cite{huang2022batterybert,xu2024atomworkbattery,dagdelen2024structured} are discussed as related work only.

\subsubsection{Warning flags}
\label{sec:warnings}

Huang and Cole attach three warning flags to records \cite{huang2020database}, and each needs a decision.

\textbf{S (series).} A record is flagged S when values occur more than three times with the same compound name, property and DOI, i.e.\ when it belongs to a data series such as capacity versus cycle number or versus current density. In the raw file 99,671 records carry S alone and 17,732 carry L and S together. The S flag is a direct function of how many sibling records a publication contributes for the same material, so it encodes publication identity rather than material properties. It is not used as a predictor. Because series members are strongly correlated with one another, they are also the records through which a random split leaks most information. We therefore report, as a sensitivity analysis, all results with the S records removed (Section~\ref{sec:series}).

\textbf{L (near limit).} A record is flagged L when its value lies in a ``near limit'' region defined by fixed value ranges (for example 0--20 and 3,000--5,000~mAh~g$^{-1}$ for capacity, 0--1 and 5--8~V for voltage, 0--20\% for Coulombic efficiency) \cite{huang2020database}. L is therefore a deterministic function of the target value. Using it as a feature would tell the model that the target lies in one of two tails, which is target leakage. L is not used as a predictor; L-flagged records are retained if they pass our own range filter, and their share of the error is reported separately (\TBD{err/L-share}).

\textbf{R (relevance).} A record is flagged R when fewer than three records were extracted from its article, a heuristic for papers that are not primarily about battery materials \cite{huang2020database}. R is constant within a publication and does not depend on the target. Because it is a paper-level constant, it is admitted only in the paper-level block (\texttt{lit\_warn\_R}, tier \texttt{t3}) together with the title keywords, and is assessed in the same way (Section~\ref{sec:controls}).

\subsubsection{Correctness sample}
\label{sec:quality}

Huang and Cole evaluated the database on about 500 records drawn by random shuffling and supplemented by hand with additional records for the rarer properties \cite{huang2020database}. Precision was 83.3\% for capacity, 79.0\% for voltage, 76.0\% for conductivity, 77.6\% for Coulombic efficiency and 75.5\% for energy (80.0\% overall); recall, estimated on 51 papers, was 59.1\% overall. A record was counted as correct when the compound name, value, property and unit all matched the source text. Errors were classified as incomplete composites (F1, e.g.\ only ``rGO'' extracted from ``N-ZnSe@rGO''), incorrect names (F2), incorrect matches between a value and a compound or property (F3, the most common type), and interdependency errors (F4). The resulting labels are stored in the \texttt{Correctness} column of the release, which we use only for analysis and never as a feature.

\subsection{Feature Engineering}
\label{sec:features}

Every feature is computed from the fields of a single record. Features are organised in blocks, identified by a column prefix, that define both the feature tiers and the groups used for SHAP aggregation (Table~\ref{tab:predictors}).

\subsubsection{Composition and Elemental-Property Descriptors}

\textbf{Composition} (\texttt{comp\_}). Atomic fractions of a predeclared list of elements, comprising every element present in at least 0.5\% of cleaned records; the list was computed once from the full corpus before any split and frozen. The block also contains the number of elements, the number of components in \texttt{Extracted\_name} (composites are parsed as several components), and indicator flags for the common carrier ions (Li, Na, K, Mg, Zn, Al).

\textbf{Elemental properties} (\texttt{mat\_}). Stoichiometry-weighted mean, minimum, maximum and range of tabulated elemental properties (atomic number, atomic mass, Pauling electronegativity, atomic radius, period, group, Mendeleev number, first ionisation energy and a common-oxidation-state proxy), computed with pymatgen \cite{ong2013pymatgen} in the manner of Magpie descriptors \cite{ward2016magpie}. Because it is computed from composition, this block adds no information that is not already implicit in the composition block, but it presents that information in a form tree models can split on.

\subsubsection{Name, Measurement and Paper-Level Features}

\textbf{Material name} (\texttt{txt\_}). Regular-expression flags on \texttt{Name} for words a materials scientist would read as architecture or morphology: carbon, graphene, CNT and rGO composites; core--shell (``@''); slash composites; doping; nano-prefixes; hollow and porous structures; films; and the words anode, cathode and electrode.

\textbf{Measurement context} (\texttt{meas\_}). Keyword flags on \texttt{Specifier} (theoretical, reversible, discharge, charge, specific, electronic/electrical, potential; the release contains only 54 distinct specifier strings, so further candidate keywords such as initial, areal, volumetric, open-circuit or cut-off never occur and were not used) and a flag for a missing specifier; current density parsed from \texttt{Info}, normalised to A~g$^{-1}$ where the unit allows and $\log_{10}$-transformed, with C-rate and cycle number where present and a missingness indicator for each; and \texttt{Type} as a categorical variable (cathode, anode, electrolyte, unknown).

\textbf{Paper-level title keywords} (\texttt{lit\_}). Flags for battery chemistry and topic words in \texttt{Title} (lithium-ion, sodium-ion, potassium, Li--S, Li--O$_2$/air, zinc, magnesium, supercapacitor, solid-state, electrolyte, anode, cathode), together with the R flag as \texttt{lit\_warn\_R}. In the released file titles are stored upper-case with spaces and stop words removed (e.g.\ \texttt{...ANODESFORLITHIUMIONBATTERIES}), so the patterns match substrings. These features are constant within a publication, which is exactly the property that makes paper-level features suspect under record-level splitting.

\textbf{Material identity} (\texttt{id\_name}). The lower-cased, whitespace-normalised \texttt{Name} string as a CatBoost categorical variable. Used only in the identity-admission control.

\textbf{Noise} (\texttt{noise\_}). 32 independent standard-normal columns, used only in the negative control.

Preprocessing was fitted on each training partition: features missing in more than 80\% of training records were dropped, except the predeclared sparse condition features (current density, C-rate, cycle number), which were kept with a missingness indicator; remaining numeric gaps were imputed with the training median and flagged.

\begin{table*}[t]
    \centering\small
    \caption{\textbf{Feature blocks and tiers.} Feature counts are after train-fitted preprocessing and can differ slightly between targets and splits; missingness indicators are counted with their base block. The main tier is $t_2$.}
    \label{tab:predictors}
    \begin{tabular}{@{}l l r p{8.2cm}@{}}
        \toprule
        Block (prefix) & Group & Features & Content \\
        \midrule
        \texttt{comp\_} & composition   & \TBDcell{feat/comp/n} & element fractions (frozen element list), number of elements and components, carrier-ion flags \\
        \texttt{mat\_}  & material      & \TBDcell{feat/mat/n} & stoichiometry-weighted mean/min/max/range of elemental properties \\
        \texttt{txt\_}  & material\_text & \TBDcell{feat/txt/n} & morphology and architecture words in the compound name \\
        \texttt{meas\_} & measurement   & \TBDcell{feat/meas/n} & specifier keywords, current density, C-rate, cycle number, \texttt{Type} \\
        \texttt{lit\_}  & literature    & \TBDcell{feat/lit/n} & title keywords and R flag (\texttt{lit\_warn\_R}; paper-level constants) \\
        \texttt{id\_name} & identity    & 1 & normalised compound name (categorical) \\
        \texttt{noise\_} & noise        & 32 & $\mathcal{N}(0,1)$ columns \\
        \midrule
        \multicolumn{4}{@{}l}{\textbf{Tiers}} \\
        $t_{\mathrm{comp}}$ & \multicolumn{3}{l}{composition only} \\
        $t_1$ & \multicolumn{3}{l}{composition + material} \\
        $t_2$ (main) & \multicolumn{3}{l}{$t_1$ + material\_text + measurement} \\
        $t_3$ & \multicolumn{3}{l}{$t_2$ + literature (title keywords and R flag)} \\
        $t_{2,\mathrm{id}}$ & \multicolumn{3}{l}{$t_2$ + \texttt{id\_name} (identity-admission control)} \\
        $t_{2,\mathrm{noise}}$ & \multicolumn{3}{l}{$t_2$ + 32 noise columns (negative control)} \\
        \bottomrule
    \end{tabular}
\end{table*}

\subsection{Data Integration and Filtering}
\label{sec:integration}

\subsubsection{Cleaning and physical-range filters}
\label{sec:cleaning}

DOIs were normalised (lower-cased; resolver prefixes removed), and records without a DOI were placed in one explicit group so that they could never be split across partitions. For each property, records were removed in the following order, and the number removed at each step is reported in Table~\ref{tab:targets}: (i) missing \texttt{Value}; (ii) energy records with \texttt{Tag} = Calculated, because they are arithmetic products of capacity and voltage records that are themselves in the database; (iii) values outside a documented physical range; and (iv) exact duplicates of (publication, property, name, value, specifier), which arise when the same sentence is parsed more than once. Records whose \texttt{Extracted\_name} could not be parsed into at least one element were kept with missing composition features.

The physical ranges (Table~\ref{tab:targets}) are deliberately wide. They are meant to remove values that cannot be correct for any battery material, not values that are merely unusual. The capacity limit of 4,500~mAh~g$^{-1}$ lies above the theoretical capacities of silicon ($\approx$4,200) and lithium metal ($\approx$3,860); the voltage limit of 6~V lies above the stability window of any electrolyte in the database; Coulombic efficiency cannot exceed 100\%; the conductivity range spans fifteen decades; and the energy limit of 3,000~Wh~kg$^{-1}$ is above any practical cell-level value reported for a material. They are tighter than the extraction limits applied by Huang and Cole (0--5,000~mAh~g$^{-1}$, 0--8~V, 0--5,000~Wh~kg$^{-1}$, $10^{-20}$--$10^{2}$~S~cm$^{-1}$) \cite{huang2020database}.

\begin{table*}[t]
    \centering\small
    \caption{\textbf{Targets, physical-range filters and records removed.} One model is trained per target. Ranges are inclusive at the upper end. Calc., Out of range and Dupl.\ are the numbers of records removed as Calculated, out of range and exact duplicates; Rec./pub.\ is the median number of records per publication. Counts are records removed at that step (Section~\ref{sec:cleaning}). Between-publication share is the one-way intraclass correlation ICC(1) of the (transformed) target by DOI.}
    \label{tab:targets}
    \setlength{\tabcolsep}{4pt}\footnotesize
    \begin{tabular}{@{}lllrrrrr@{}}
        \toprule
        Target & Unit (model scale) & Kept range & Calc. & Out of range & Dupl. & Rec./pub. & ICC(1) \\
        \midrule
        Capacity              & mAh g$^{-1}$ & $0 < v \le 4500$ & --- & \TBDcell{pop/capacity/removed_bounds} & \TBDcell{pop/capacity/removed_duplicates} & \TBDcell{pop/capacity/records_per_pub_median} & \TBDcell{pop/capacity/icc1} \\
        Voltage               & V & $0 < v \le 6$    & --- & \TBDcell{pop/voltage/removed_bounds} & \TBDcell{pop/voltage/removed_duplicates} & \TBDcell{pop/voltage/records_per_pub_median} & \TBDcell{pop/voltage/icc1} \\
        Conductivity          & $\log_{10}$ S cm$^{-1}$ & $10^{-12} \le v \le 10^{3}$ & --- & \TBDcell{pop/conductivity/removed_bounds} & \TBDcell{pop/conductivity/removed_duplicates} & \TBDcell{pop/conductivity/records_per_pub_median} & \TBDcell{pop/conductivity/icc1} \\
        Coulombic efficiency  & \% & $0 < v \le 100$  & --- & \TBDcell{pop/coulombic_efficiency/removed_bounds} & \TBDcell{pop/coulombic_efficiency/removed_duplicates} & \TBDcell{pop/coulombic_efficiency/records_per_pub_median} & \TBDcell{pop/coulombic_efficiency/icc1} \\
        Energy                & Wh kg$^{-1}$ & $0 < v \le 3000$ & \TBDcell{pop/energy/removed_calculated} & \TBDcell{pop/energy/removed_bounds} & \TBDcell{pop/energy/removed_duplicates} & \TBDcell{pop/energy/records_per_pub_median} & \TBDcell{pop/energy/icc1} \\
        \bottomrule
    \end{tabular}
\end{table*}

\subsubsection{Excluded columns}
\label{sec:forbidden}

Table~S2 lists every raw column that is never used directly as a predictor and the reason. The list is enforced in code: the feature-frame constructor raises an error if any of these columns reaches a model. Several columns enter only through a derived, documented block; the raw text itself is excluded because a free-text field can act as a record or publication identifier.



\subsection{Machine Learning Models and Evaluation}

\subsubsection{Model}
\label{sec:models}

The model is an equal-weight ensemble of CatBoost \cite{prokhorenkova2018catboost} and XGBoost \cite{chen2016xgboost} (Table~S3). In each split a CatBoost configuration is selected from a predeclared six-point grid of tree depth and L2 regularisation by validation RMSE, refitted with three seeds, and the three predictions averaged. XGBoost is early-stopped on the same validation set, with categorical variables ordinally encoded on the training partition. The ensemble prediction is $\hat y = \tfrac12\,\overline{\hat y}_{\mathrm{CatBoost}} + \tfrac12\,\hat y_{\mathrm{XGBoost}}$. Both learners receive the full list of categorical columns present in the frame (\texttt{meas\_type}, and \texttt{id\_name} in $t_{2,\mathrm{id}}$). No hyperparameter was tuned on test data, and the same settings are used for all targets, tiers and protocols. No training partition is subsampled: every target, tier and protocol is trained on its full training split (up to about 98,000 records for capacity).



\subsubsection{Evaluation protocols}
\label{sec:protocols}

Three split protocols were applied to every target with identical model, features and preprocessing (Table~S4). Splits use scikit-learn \cite{pedregosa2011sklearn}.

\textbf{Record-level random.} Five repeated \texttt{ShuffleSplit} holdouts (test fraction 0.2, seed 42), with a validation set of 15\% of the remaining records for early stopping and model selection.

\textbf{DOI-grouped.} Five repeated \texttt{GroupShuffleSplit} holdouts with the same fractions and seed, grouped by normalised DOI; the validation set is drawn by publication as well, so training, validation and test publications are disjoint. A disjoint five-fold grouped variant, in which the test folds partition the publications, is used to check that the repeated-holdout estimates are not affected by overlap between test sets.

\textbf{Temporal.} Publications are dated by the earliest parsed date of their records, and records from undated publications are excluded from this protocol only. Publications are ranked by date, and three forward splits test on the publications in the 55--70th, 70--85th and 85--100th date percentiles, respectively; training uses only earlier publications, with the most recent 15\% of those as the validation set. This protocol asks how well a model trained on the literature up to some date predicts what was published afterwards, which includes new chemistries that were absent from training \cite{sheridan2013timesplit}.

\textbf{Composition-grouped.} Records are grouped by their parsed composition (atomic fractions of the frozen element list, rounded to two decimals; 7,708 groups for capacity, the largest holding 4\% of records) and split with the same \texttt{GroupShuffleSplit} settings, so that no composition occurs in more than one partition. Publications may now be shared between partitions. This protocol approximates the leave-cluster-out evaluations used to measure extrapolation to new materials \cite{meredig2018cluster,witman2025matfold}.

For each protocol we report the share of test records whose publication also contributes training records (by construction zero for the DOI-grouped and temporal splits) and the share whose normalised compound name occurs in training. The second quantity matters because DOI-grouped splitting removes unseen \emph{publications} but not unseen \emph{materials}: LiFePO$_4$ or graphite appear in thousands of papers, and a DOI-grouped test set still contains many materials that the model has seen.



\textbf{Metrics.} For each test partition we compute $R^2$, MAE and RMSE between the reported values $y_i$ and predictions $\hat y_i$,
\begin{equation}
    R^2 = 1-\frac{\sum_{i}(y_i-\hat{y}_i)^2}{\sum_{i}(y_i-\bar{y})^2},\quad
    \mathrm{MAE} = \frac{1}{N}\sum_{i}|y_i-\hat{y}_i|,\quad
    \mathrm{RMSE} = \sqrt{\frac{1}{N}\sum_{i}(y_i-\hat{y}_i)^2},
\end{equation}
with $\bar y$ the mean of the test partition. Results are the mean $\pm$ sample standard deviation over splits. Because $R^2$ depends on the variance of the test set, which differs between protocols, MAE and RMSE are the primary comparison across protocols and $R^2$ the primary comparison across targets. As a reference we also report a constant-mean baseline and a per-specifier median baseline, which predicts the training median of records with the same specifier keyword set.

\subsubsection{Controls}
\label{sec:controls}

\textbf{Negative noise control.} The $t_{2,\mathrm{noise}}$ tier adds 32 standard-normal columns. A correctly regularised model should give them a negligible share of the SHAP attribution and should not improve on $t_2$; a substantial share would indicate that the model is fitting noise.

\textbf{Paper-level features.} The $t_3$ tier adds title keywords and the R flag. Because these features are constant within a publication, they can identify groups of publications. We compare $t_3$ with $t_2$ under both random and DOI-grouped splits: a gain under random splitting that vanishes under DOI grouping indicates that the features help by identifying publications rather than by describing materials.

\textbf{Identity admission.} The $t_{2,\mathrm{id}}$ tier adds the compound name as a categorical variable. With CatBoost's ordered target statistics this lets the model use the mean target of each named material. The comparison tests how much of the predictive signal is attached to the name rather than to the parsed composition. We report it as a control and do not include it in the main model.

\textbf{Series sensitivity.} All protocols are repeated on the data with S-flagged records removed (Section~\ref{sec:warnings}).

For tier comparisons we report paired differences over the same splits, with the publication-cluster bootstrap 95\% interval of $\Delta R^2$ on the pooled test predictions.

\subsubsection{Uncertainty}
\label{sec:uncertainty_methods}

Besides the spread over repeated holdouts, we use the disagreement between the four ensemble members (three CatBoost seeds and XGBoost) as a per-record epistemic proxy, $s_i = \mathrm{sd}(\hat y_i^{(1)},\ldots,\hat y_i^{(4)})$. We evaluate (i) the Spearman rank correlation between $s_i$ and $|y_i - \hat y_i|$ in each test partition; (ii) whether $s_i$ is larger for DOI-grouped than for record-level test records; and (iii) the error of the records in the top and bottom deciles of $s_i$. This is a cheap proxy, not a calibrated predictive interval: three of the four members share data and hyperparameters, so their spread mainly reflects variance between models and not noise in the measurements \cite{tran2020uq}.

\subsubsection{Leakage-channel analyses}
\label{sec:channels_methods}

Three further analyses use the stored test predictions and require no refitting. (i) \emph{Same-publication reference.} Under record-level splitting, each test record is predicted by the mean of the training records from its own publication (the training mean when there are none). This model knows nothing about materials and measures how much of the target can be recovered from publication identity alone. (ii) \emph{Values recurring across publications.} A test record is called \emph{repeated} when the same normalised compound name and the same value occur in a training record from any publication; under DOI-grouped splitting this is necessarily a different publication. We report the share of repeated test records and $R^2$ separately for repeated and non-repeated records. Recurring values arise from reported theoretical capacities, benchmark plateaus, and values quoted from earlier work. (iii) \emph{Theoretical capacities.} Records whose specifier contains ``theoretical'' are functions of the composition by definition; we report capacity metrics with these records excluded. Finally, to test whether publication identity is recoverable from the descriptors, we repeat the bibliographic-fingerprint test of Jablonka \cite{jablonka2026clever} on our splits (Section~\ref{sec:previous}).

\subsection{Feature Importance and Model Interpretation}
\label{sec:shap_methods}

Exact TreeSHAP values \cite{shappaper,lundberg2020treeshap} were computed for each CatBoost seed model and the XGBoost model on the test records of every DOI-grouped split. Because the ensemble is linear in its members, the ensemble attribution is $\phi_{\mathrm{ens}} = \tfrac12\overline{\phi}_{\mathrm{CatBoost}} + \tfrac12\phi_{\mathrm{XGBoost}}$. The share of a group is its summed mean $|$SHAP$|$ divided by the total over all predictors; missingness indicators are assigned to the group of their base feature. Shares are reported as mean $\pm$ sample standard deviation over splits. SHAP values describe the fitted model and do not establish causal effects.

% ---------------- RESULTS AND DISCUSSION ----------------
\section{Results and Discussion}

\subsection{Performance of the final ensemble under record-level, DOI-grouped and temporal evaluation}
\label{sec:splits}

Table~\ref{tab:dataset} and Table~\ref{tab:targets} summarise the modelling sets. After cleaning, the five targets contain \TBD{pop/total/records} records from \TBD{pop/total/publications} publications. The range filters removed \TBD{pop/all/removed_bounds} records in total, most of them \TBD{pop/removed_bounds/which-target}. Publications contribute a median of \TBD{pop/capacity/records_per_pub_median} capacity records but only \TBD{pop/conductivity/records_per_pub_median} conductivity record, so the potential for within-publication leakage differs strongly between targets. The between-publication share of target variance, ICC(1), ranges from \TBD{pop/icc1/min} (\TBD{pop/icc1/argmin}) to \TBD{pop/icc1/max} (\TBD{pop/icc1/argmax}).

\begin{figure}[htbp]
    \centering
    \FIGTBD{figures/fig_target_distributions.png}{Five panels: distribution of each (transformed) target after cleaning; S-flagged records overlaid.}
    \caption{\textbf{Distributions of the five targets after cleaning.} Conductivity is shown as $\log_{10}(\sigma/\mathrm{S\,cm^{-1}})$. The shaded histogram shows the records flagged as series members (S).}
    \label{fig:dist}
\end{figure}

Table~\ref{tab:split} and Fig.~\ref{fig:split} report the main model ($t_2$) for each target under the three protocols. For capacity, $R^2$ fell from \TBD{m/capacity/random/t2/r2} under record-level splitting to \TBD{m/capacity/doi/t2/r2} under DOI-grouped splitting and \TBD{m/capacity/temporal/t2/r2} under temporal splitting, and MAE rose from \TBD{m/capacity/random/t2/mae} to \TBD{m/capacity/doi/t2/mae}~mAh~g$^{-1}$. The other targets behave in the same way. DOI grouping lowered $R^2$ from \TBD{m/voltage/random/t2/r2} to \TBD{m/voltage/doi/t2/r2} for voltage, from \TBD{m/conductivity/random/t2/r2} to \TBD{m/conductivity/doi/t2/r2} for conductivity, from \TBD{m/coulombic_efficiency/random/t2/r2} to \TBD{m/coulombic_efficiency/doi/t2/r2} for Coulombic efficiency and from \TBD{m/energy/random/t2/r2} to \TBD{m/energy/doi/t2/r2} for energy. The gaps (0.02--0.08) are small given the strong between-publication structure, and they are not ordered by ICC(1): energy, with the highest between-publication share (0.84), has the largest gap, but voltage (0.67) has one of the smallest. A large between-publication share is therefore necessary but not sufficient for a large record-level advantage; the model must also be able to recognise the publication from its features (Section~\ref{sec:channels}). Every model beats the constant-mean baseline under every protocol, but Coulombic efficiency only marginally ($R^2 \le 0.1$). The per-specifier median baseline was worse than the constant mean for every target (DOI-grouped capacity RMSE \TBD{m/capacity/doi/specbase/rmse} against \TBD{m/capacity/doi/baseline/rmse}~mAh~g$^{-1}$), so the specifier alone explains little of the variance; its value lies in interaction with the composition..

Under record-level splitting, \TBD{overlap/random/doi_seen_share} of test records came from publications that also contributed training records. Under DOI-grouped splitting this share is zero, but \TBD{overlap/doi/name_seen_share} of test records had a compound name that occurs in training, compared with \TBD{overlap/random/name_seen_share} under random splitting and \TBD{overlap/temporal/name_seen_share} under temporal splitting. The DOI-grouped result is therefore an estimate of performance on new papers about largely familiar materials, and the temporal result is closer to performance on new materials. Temporal splitting was the strictest of the three publication-level protocols for every target: relative to DOI grouping, capacity lost a further \TBD{d/capacity/temporal-doi/r2} in $R^2$, and the shortfall was largest for energy. Later publications report more materials that are absent from earlier ones (\TBD{overlap/temporal/name_seen_share} of temporal test records had a known name, against \TBD{overlap/doi/name_seen_share} under DOI grouping), so the temporal estimate is closer to prospective use. Holding out whole compositions was stricter than DOI grouping for every target and stricter than the temporal split for voltage, conductivity and energy, with a large spread between splits: capacity reached \TBD{m/capacity/material/t2/r2}, voltage \TBD{m/voltage/material/t2/r2} and conductivity \TBD{m/conductivity/material/t2/r2} (Table~\ref{tab:channels}). The disjoint five-fold grouped check reproduced the repeated-holdout DOI estimates (capacity \TBD{m/capacity/kfold/t2/r2}).

\begin{table*}[t]
    \centering\small
    \caption{\textbf{Main model ($t_2$) under record-level random, DOI-grouped and temporal evaluation.} Mean $\pm$ sample standard deviation over splits (5, 5 and 3). $N$ is the mean number of test records per split. Units of MAE and RMSE: mAh~g$^{-1}$ (capacity), V (voltage), decades of S~cm$^{-1}$ (conductivity), percentage points (Coulombic efficiency), Wh~kg$^{-1}$ (energy). The last column is the constant-mean baseline RMSE on the DOI-grouped splits.}
    \label{tab:split}
    \begin{tabular}{@{}llrcccc@{}}
        \toprule
        Target & Protocol & $N$ & $R^2$ & MAE & RMSE & Baseline RMSE \\
        \midrule
        Capacity & Random   & \TBDcell{m/capacity/random/t2/n} & \TBDcell{m/capacity/random/t2/r2} & \TBDcell{m/capacity/random/t2/mae} & \TBDcell{m/capacity/random/t2/rmse} & \\
                 & DOI      & \TBDcell{m/capacity/doi/t2/n} & \TBDcell{m/capacity/doi/t2/r2} & \TBDcell{m/capacity/doi/t2/mae} & \TBDcell{m/capacity/doi/t2/rmse} & \TBDcell{m/capacity/doi/baseline/rmse} \\
                 & Temporal & \TBDcell{m/capacity/temporal/t2/n} & \TBDcell{m/capacity/temporal/t2/r2} & \TBDcell{m/capacity/temporal/t2/mae} & \TBDcell{m/capacity/temporal/t2/rmse} & \\
        \midrule
        Voltage  & Random   & \TBDcell{m/voltage/random/t2/n} & \TBDcell{m/voltage/random/t2/r2} & \TBDcell{m/voltage/random/t2/mae} & \TBDcell{m/voltage/random/t2/rmse} & \\
                 & DOI      & \TBDcell{m/voltage/doi/t2/n} & \TBDcell{m/voltage/doi/t2/r2} & \TBDcell{m/voltage/doi/t2/mae} & \TBDcell{m/voltage/doi/t2/rmse} & \TBDcell{m/voltage/doi/baseline/rmse} \\
                 & Temporal & \TBDcell{m/voltage/temporal/t2/n} & \TBDcell{m/voltage/temporal/t2/r2} & \TBDcell{m/voltage/temporal/t2/mae} & \TBDcell{m/voltage/temporal/t2/rmse} & \\
        \midrule
        Conductivity ($\log_{10}$) & Random & \TBDcell{m/conductivity/random/t2/n} & \TBDcell{m/conductivity/random/t2/r2} & \TBDcell{m/conductivity/random/t2/mae} & \TBDcell{m/conductivity/random/t2/rmse} & \\
                 & DOI      & \TBDcell{m/conductivity/doi/t2/n} & \TBDcell{m/conductivity/doi/t2/r2} & \TBDcell{m/conductivity/doi/t2/mae} & \TBDcell{m/conductivity/doi/t2/rmse} & \TBDcell{m/conductivity/doi/baseline/rmse} \\
                 & Temporal & \TBDcell{m/conductivity/temporal/t2/n} & \TBDcell{m/conductivity/temporal/t2/r2} & \TBDcell{m/conductivity/temporal/t2/mae} & \TBDcell{m/conductivity/temporal/t2/rmse} & \\
        \midrule
        Coulombic eff. & Random & \TBDcell{m/coulombic_efficiency/random/t2/n} & \TBDcell{m/coulombic_efficiency/random/t2/r2} & \TBDcell{m/coulombic_efficiency/random/t2/mae} & \TBDcell{m/coulombic_efficiency/random/t2/rmse} & \\
                 & DOI      & \TBDcell{m/coulombic_efficiency/doi/t2/n} & \TBDcell{m/coulombic_efficiency/doi/t2/r2} & \TBDcell{m/coulombic_efficiency/doi/t2/mae} & \TBDcell{m/coulombic_efficiency/doi/t2/rmse} & \TBDcell{m/coulombic_efficiency/doi/baseline/rmse} \\
                 & Temporal & \TBDcell{m/coulombic_efficiency/temporal/t2/n} & \TBDcell{m/coulombic_efficiency/temporal/t2/r2} & \TBDcell{m/coulombic_efficiency/temporal/t2/mae} & \TBDcell{m/coulombic_efficiency/temporal/t2/rmse} & \\
        \midrule
        Energy   & Random   & \TBDcell{m/energy/random/t2/n} & \TBDcell{m/energy/random/t2/r2} & \TBDcell{m/energy/random/t2/mae} & \TBDcell{m/energy/random/t2/rmse} & \\
                 & DOI      & \TBDcell{m/energy/doi/t2/n} & \TBDcell{m/energy/doi/t2/r2} & \TBDcell{m/energy/doi/t2/mae} & \TBDcell{m/energy/doi/t2/rmse} & \TBDcell{m/energy/doi/baseline/rmse} \\
                 & Temporal & \TBDcell{m/energy/temporal/t2/n} & \TBDcell{m/energy/temporal/t2/r2} & \TBDcell{m/energy/temporal/t2/mae} & \TBDcell{m/energy/temporal/t2/rmse} & \\
        \bottomrule
    \end{tabular}
\end{table*}

\begin{figure}[htbp]
    \centering
    \FIGTBD{figures/fig_split_by_target.png}{Grouped bars: R2 (mean $\pm$ sd) per target under random (hatched), DOI-grouped (solid) and temporal (dotted).}
    \caption{\textbf{Effect of the split protocol for each target.} $R^2$ of the $t_2$ ensemble under record-level random, DOI-grouped and temporal evaluation; bars and error bars are mean and sample standard deviation over splits.}
    \label{fig:split}
\end{figure}

\subsection{Why the DOI-grouped gap is small: leakage channels}
\label{sec:channels}

The small gap could mean that records from one publication are nearly independent, but Table~\ref{tab:targets} shows the opposite. Under record-level splitting, predicting each test record by the mean of the training records from its own publication, with no information about the material, gives $R^2$ of \TBD{oracle/capacity/samepaper_r2} for capacity, \TBD{oracle/voltage/samepaper_r2} for voltage and \TBD{oracle/energy/samepaper_r2} for energy, a large fraction of what the full model achieves (and nearly all of it for voltage). Publication identity is therefore a strong predictor. The model does not exploit it because its features cannot recover it: from the composition descriptors, the journal of a record is predicted with a micro-F1 of \TBD{ch/capacity/random/journal_f1} when publications are shared between training and test but only \TBD{ch/capacity/doi/journal_f1} when they are not, the publication year barely improves on the training mean, and a model given only the predicted journal and year has $R^2 \approx 0$ under both splits (Section~\ref{sec:previous}). In databases whose records carry process settings, device architectures or characteristic patterns of missing fields, such combinations can be specific to a laboratory and let a record-level split recognise the source; the battery records carry almost none of them.

A second channel is not blocked by DOI grouping. Under DOI-grouped splitting, \TBD{rep/capacity/doi/share} of capacity test records, \TBD{rep/voltage/doi/share} of voltage records and \TBD{rep/conductivity/doi/share} of conductivity records repeat a (compound name, value) pair that another publication reported in the training set (Table~\ref{tab:channels}). These records are predicted far better than the rest: capacity $R^2$ was \TBD{rep/capacity/doi/r2_rep} on repeated and \TBD{rep/capacity/doi/r2_new} on non-repeated records, and voltage \TBD{rep/voltage/doi/r2_rep} against \TBD{rep/voltage/doi/r2_new}. The repeated share and its advantage are almost the same under record-level splitting (capacity \TBD{rep/capacity/random/share}, $R^2$ \TBD{rep/capacity/random/r2_rep}), which is why the two protocols give similar overall scores: both reward recall of values that recur across the literature. For conductivity and Coulombic efficiency, records without such a precedent are predicted hardly better than their mean ($R^2$ \TBD{rep/conductivity/doi/r2_new} and \TBD{rep/coulombic_efficiency/doi/r2_new}). Theoretical capacities, which are a deterministic function of composition, make up \TBD{theo/capacity/share} of DOI-grouped capacity test records; excluding them lowers capacity $R^2$ from \TBD{m/capacity/doi/t2/r2} to \TBD{notheo/capacity/doi/r2}.

\begin{table*}[t]
    \centering\small
    \caption{\textbf{Leakage channels by target} ($t_2$ ensemble; mean over splits). Same-pub.: $R^2$ of the same-publication mean under record-level splitting. Repeated: share of DOI-grouped test records whose (name, value) pair occurs in a training record of another publication, with $R^2$ on repeated and non-repeated records. Comp.-grouped: $R^2$ under composition-grouped splitting.}
    \label{tab:channels}
    \begin{tabular}{@{}lccccc@{}}
        \toprule
        Target & Same-pub.\ $R^2$ & Repeated share (DOI) & $R^2$ repeated & $R^2$ non-repeated & Comp.-grouped $R^2$ \\
        \midrule
        Capacity       & \TBDcell{oracle/capacity/samepaper_r2} & \TBDcell{rep/capacity/doi/share} & \TBDcell{rep/capacity/doi/r2_rep} & \TBDcell{rep/capacity/doi/r2_new} & \TBDcell{m/capacity/material/t2/r2} \\
        Voltage        & \TBDcell{oracle/voltage/samepaper_r2} & \TBDcell{rep/voltage/doi/share} & \TBDcell{rep/voltage/doi/r2_rep} & \TBDcell{rep/voltage/doi/r2_new} & \TBDcell{m/voltage/material/t2/r2} \\
        Conductivity   & \TBDcell{oracle/conductivity/samepaper_r2} & \TBDcell{rep/conductivity/doi/share} & \TBDcell{rep/conductivity/doi/r2_rep} & \TBDcell{rep/conductivity/doi/r2_new} & \TBDcell{m/conductivity/material/t2/r2} \\
        Coulombic eff. & \TBDcell{oracle/coulombic_efficiency/samepaper_r2} & \TBDcell{rep/coulombic_efficiency/doi/share} & \TBDcell{rep/coulombic_efficiency/doi/r2_rep} & \TBDcell{rep/coulombic_efficiency/doi/r2_new} & \TBDcell{m/coulombic_efficiency/material/t2/r2} \\
        Energy         & \TBDcell{oracle/energy/samepaper_r2} & \TBDcell{rep/energy/doi/share} & \TBDcell{rep/energy/doi/r2_rep} & \TBDcell{rep/energy/doi/r2_new} & \TBDcell{m/energy/material/t2/r2} \\
        \bottomrule
    \end{tabular}
\end{table*}

\subsection{Effect of the paper-level (title-keyword) features and controls}
\label{sec:tiers}

Table~\ref{tab:tiers} reports DOI-grouped $R^2$ for every tier and target. Adding elemental-property statistics to composition ($t_{\mathrm{comp}} \rightarrow t_1$) changed $R^2$ by \TBD{d/capacity/doi/t1-comp/r2} for capacity, and adding the name and measurement-context blocks ($t_1 \rightarrow t_2$) by \TBD{d/capacity/doi/t2-t1/r2}. Elemental-property statistics add almost nothing to the element fractions, as expected for a block that is a function of composition. The name and measurement-context blocks add most for capacity, where the specifier separates theoretical from measured values and current density and cycle number carry the measurement conditions, and little for Coulombic efficiency..

\textbf{Noise control.} The 32 noise columns changed DOI-grouped $R^2$ by \TBD{d/all/doi/noise-t2/r2-range} across targets and received \TBD{shap/noise/share} of the total SHAP attribution (Table~\ref{tab:shap}). The models do not benefit from noise, but for the weak conductivity and Coulombic-efficiency models the noise columns receive about a quarter of the attribution: when the signal is weak, TreeSHAP spreads attribution over whatever the trees split on, and the feature rankings for these two targets should not be interpreted.

\textbf{Paper-level title keywords.} Adding title keywords and the R flag ($t_3 - t_2$) changed $R^2$ by \TBD{d/capacity/random/t3-t2/r2} under random splitting and by \TBD{d/capacity/doi/t3-t2/r2} \TBD{ci/capacity/doi/t3-t2/r2} under DOI-grouped splitting for capacity (Table~S5). Unlike the paper-level features that are usually suspected of identifying publications, the title keywords help almost as much under DOI grouping as under random splitting for capacity and voltage, somewhat less for energy and Coulombic efficiency, and not significantly for conductivity; the gain is largest for energy ($\Delta R^2$ \TBD{d/energy/doi/t3-t2/r2} \TBD{ci/energy/doi/t3-t2/r2}). They therefore carry chemistry rather than identity: the formula of the active material does not say whether it was tested in a lithium-ion, sodium-ion, Li--S or Li--O$_2$ cell, and the title does. Because the information is constant within a publication, a model that relies on it is still a model of how a paper's system is described, and we keep it out of the main tier..

\textbf{Identity admission.} Adding the compound name as a categorical variable ($t_{2,\mathrm{id}} - t_2$) changed DOI-grouped $R^2$ by \TBD{d/capacity/doi/t2_id-t2/r2} for capacity and by \TBD{d/all/doi/t2_id-t2/r2-range} across targets. Under DOI grouping the effect is within $\pm0.01$ for every target (positive only for energy, \TBD{d/energy/doi/t2_id-t2/r2} \TBD{ci/energy/doi/t2_id-t2/r2}); under random splitting the name added slightly more (capacity \TBD{d/capacity/random/t2_id-t2/r2}), a mild identity effect that largely disappears when publications are separated. Because the composition block already nearly determines the material, the name adds little beyond it.

\begin{table*}[t]
    \centering\small
    \caption{\textbf{DOI-grouped $R^2$ by feature tier} (mean $\pm$ sample standard deviation over five splits; same splits for all tiers).}
    \label{tab:tiers}
    \begin{tabular}{@{}lccccc@{}}
        \toprule
        Tier & Capacity & Voltage & Conductivity & Coulombic eff. & Energy \\
        \midrule
        $t_{\mathrm{comp}}$     & \TBDcell{m/capacity/doi/comp/r2} & \TBDcell{m/voltage/doi/comp/r2} & \TBDcell{m/conductivity/doi/comp/r2} & \TBDcell{m/coulombic_efficiency/doi/comp/r2} & \TBDcell{m/energy/doi/comp/r2} \\
        $t_1$                   & \TBDcell{m/capacity/doi/t1/r2} & \TBDcell{m/voltage/doi/t1/r2} & \TBDcell{m/conductivity/doi/t1/r2} & \TBDcell{m/coulombic_efficiency/doi/t1/r2} & \TBDcell{m/energy/doi/t1/r2} \\
        $t_2$ (main)            & \TBDcell{m/capacity/doi/t2/r2} & \TBDcell{m/voltage/doi/t2/r2} & \TBDcell{m/conductivity/doi/t2/r2} & \TBDcell{m/coulombic_efficiency/doi/t2/r2} & \TBDcell{m/energy/doi/t2/r2} \\
        $t_3$                   & \TBDcell{m/capacity/doi/t3/r2} & \TBDcell{m/voltage/doi/t3/r2} & \TBDcell{m/conductivity/doi/t3/r2} & \TBDcell{m/coulombic_efficiency/doi/t3/r2} & \TBDcell{m/energy/doi/t3/r2} \\
        $t_{2,\mathrm{id}}$     & \TBDcell{m/capacity/doi/t2_id/r2} & \TBDcell{m/voltage/doi/t2_id/r2} & \TBDcell{m/conductivity/doi/t2_id/r2} & \TBDcell{m/coulombic_efficiency/doi/t2_id/r2} & \TBDcell{m/energy/doi/t2_id/r2} \\
        $t_{2,\mathrm{noise}}$  & \TBDcell{m/capacity/doi/t2_noise/r2} & \TBDcell{m/voltage/doi/t2_noise/r2} & \TBDcell{m/conductivity/doi/t2_noise/r2} & \TBDcell{m/coulombic_efficiency/doi/t2_noise/r2} & \TBDcell{m/energy/doi/t2_noise/r2} \\
        \bottomrule
    \end{tabular}
\end{table*}



\subsection{Prediction distribution and errors}
\label{sec:prediction_diagnostics}

\textbf{Member spread and absolute error.} The Spearman correlation between member spread and absolute error was \TBD{unc/capacity/doi/spearman} for capacity under DOI-grouped evaluation and \TBD{unc/capacity/random/spearman} under random evaluation (Table~S6). The mean spread for DOI-grouped test records was \TBD{unc/capacity/spread_ratio_doi_over_random} times that for record-level test records. Records in the top decile of spread had an MAE of \TBD{unc/capacity/doi/mae_top_decile}, compared with \TBD{unc/capacity/doi/mae_bottom_decile} in the bottom decile. The spread is informative only for capacity and energy ($\rho \approx 0.35$); for voltage, conductivity and Coulombic efficiency $\rho \le 0.16$. It rises only slightly for publications not seen in training (ratio 1.03--1.17), much less than the error does. Four members that share data and hyperparameters are not a substitute for a calibrated interval, and we use the spread only to rank records for inspection.



\begin{figure}[htbp]
    \centering
    \FIGTBD{figures/fig_uncertainty.png}{(a) |error| vs member spread (hexbin, capacity, DOI-grouped); (b) MAE by spread decile, random vs DOI-grouped, all targets.}
    \caption{\textbf{Ensemble spread and error.} (a) Absolute error versus member spread for capacity on DOI-grouped test records. (b) MAE by decile of spread for each target under random and DOI-grouped evaluation.}
    \label{fig:unc}
\end{figure}

\textbf{High-value errors.} Predictions regress towards the mean. Under DOI-grouped evaluation the decile of highest reported capacities (mean \TBD{hv/capacity/top/ymean}~mAh~g$^{-1}$) was under-predicted by \TBD{hv/capacity/top/bias}~mAh~g$^{-1}$ on average and the lowest decile (mean \TBD{hv/capacity/bottom/ymean}) over-predicted by \TBD{hv/capacity/bottom/bias}; the same pattern holds for voltage (top decile \TBD{hv/voltage/top/ymean}~V, bias \TBD{hv/voltage/top/bias}~V) and energy. Many of the extreme values are near the L-flag limits and include extraction errors, so part of this bias is the correct response of a regularised model to label noise; the model should not be used to rank candidate materials by extreme predicted capacity..

\textbf{Within-paper series.}\label{sec:series} Removing S-flagged records reduced the capacity set to \TBD{ser/capacity/records} records. The record-level minus DOI-grouped gap in $R^2$ was \TBD{ser/capacity/gap-with-S} with series records and \TBD{ser/capacity/gap-without-S} without them (Table~\ref{tab:series}). Series members are therefore the main channel through which record-level splitting rewards within-publication similarity: without them the record-level advantage shrinks for every target except Coulombic efficiency, and both protocols score higher, because series members are the records whose values depend most on measurement conditions (current density, cycle number) that the features describe only partly. Removing S also changes the population, so the absolute values with and without S are not a like-for-like comparison of model quality.

\begin{table}[htbp]
    \centering\small
    \caption{\textbf{Sensitivity to within-paper series.} $R^2$ of the $t_2$ ensemble with all records and with S-flagged records removed.}
    \label{tab:series}
    \begin{tabular}{@{}lcccc@{}}
        \toprule
        & \multicolumn{2}{c}{All records} & \multicolumn{2}{c}{Without S} \\
        \cmidrule(lr){2-3}\cmidrule(l){4-5}
        Target & Random & DOI & Random & DOI \\
        \midrule
        Capacity       & \TBDcell{m/capacity/random/t2/r2} & \TBDcell{m/capacity/doi/t2/r2} & \TBDcell{ser/capacity/random/r2} & \TBDcell{ser/capacity/doi/r2} \\
        Voltage        & \TBDcell{m/voltage/random/t2/r2} & \TBDcell{m/voltage/doi/t2/r2} & \TBDcell{ser/voltage/random/r2} & \TBDcell{ser/voltage/doi/r2} \\
        Conductivity   & \TBDcell{m/conductivity/random/t2/r2} & \TBDcell{m/conductivity/doi/t2/r2} & \TBDcell{ser/conductivity/random/r2} & \TBDcell{ser/conductivity/doi/r2} \\
        Coulombic eff. & \TBDcell{m/coulombic_efficiency/random/t2/r2} & \TBDcell{m/coulombic_efficiency/doi/t2/r2} & \TBDcell{ser/coulombic_efficiency/random/r2} & \TBDcell{ser/coulombic_efficiency/doi/r2} \\
        Energy         & \TBDcell{m/energy/random/t2/r2} & \TBDcell{m/energy/doi/t2/r2} & \TBDcell{ser/energy/random/r2} & \TBDcell{ser/energy/doi/r2} \\
        \bottomrule
    \end{tabular}
\end{table}

\textbf{Data quality of the extracted records.} Four consequences of the extraction errors described in Section~\ref{sec:quality} follow. First, with roughly one record in five wrong in at least one field, the achievable $R^2$ is bounded well below one even for a perfect model, and errors in the value or the property assignment act as label noise that no feature can explain. Second, F1 and F2 errors corrupt the composition features themselves; the raw file contains, for example, a record named ``Carbon'' whose parsed composition is C$_{12}$H$_{22}$O$_{11}$ (sucrose, presumably the carbon source in that paper). Third, because the evaluation sample was augmented for rare properties, it is not a simple random sample, and the precision values should be taken as approximate. Fourth, the Correctness sample is too small to estimate model error separately for correct and incorrect records with any precision; we report descriptively that the DOI-grouped absolute error of the \TBD{corr/n_test} labelled records that fell in test partitions was \TBD{corr/mae_T} for records marked correct and \TBD{corr/mae_F} for those marked incorrect. Our physical-range filter removes only the grossest errors; values that are wrong but plausible remain.

\subsection{Interpretation of model predictions}
\label{sec:shap}

Table~\ref{tab:shap} and Fig.~\ref{fig:shap} give the share of mean $|$SHAP$|$ by feature group for the $t_2$ ensembles under DOI-grouped evaluation (and for the noise columns in $t_{2,\mathrm{noise}}$). For capacity, measurement-context features accounted for \TBD{shap/capacity/measurement/share}, composition for \TBD{shap/capacity/composition/share}, elemental properties for \TBD{shap/capacity/material/share} and name flags for \TBD{shap/capacity/material_text/share}. The highest-ranked individual feature was \TBD{shapf/capacity/rank1/feature} (\TBD{shapf/capacity/rank1/mean_abs_shap}~mAh~g$^{-1}$). Elemental-property statistics carry the largest share for every target (50--70\%), composition fractions most of the remainder, and the name flags less than 5\%. Measurement context matters mainly for capacity (\TBD{shap/capacity/measurement/share}), where the specifier and current density distinguish theoretical from measured and low-rate from high-rate values, and hardly at all for Coulombic efficiency. The leading individual features are chemically plausible: silicon fraction and ionisation-energy and group statistics for capacity, Mendeleev-number and electronegativity ranges and lithium fraction for voltage, and lithium content for energy, which is dominated by Li--S and Li--O$_2$ systems at the high end. For conductivity and Coulombic efficiency the noise control shows that rankings are unreliable (Section~\ref{sec:tiers}).

Attributions to composition and elemental-property features should be read with care. Elemental statistics are functions of composition and can act as a coordinate system for recognising familiar materials; a high attribution does not mean the model has learned a physical relationship between, for example, electronegativity and voltage.

\begin{table}[htbp]
    \centering\small
    \caption{\textbf{Share of mean $|$SHAP$|$ by feature group} ($t_2$ ensembles, DOI-grouped test records; mean over five splits; noise share from $t_{2,\mathrm{noise}}$).}
    \label{tab:shap}
    \begin{tabular}{@{}lccccc@{}}
        \toprule
        Group & Capacity & Voltage & Conductivity & Coulombic eff. & Energy \\
        \midrule
        Composition    & \TBDcell{shap/capacity/composition/share} & \TBDcell{shap/voltage/composition/share} & \TBDcell{shap/conductivity/composition/share} & \TBDcell{shap/coulombic_efficiency/composition/share} & \TBDcell{shap/energy/composition/share} \\
        Material       & \TBDcell{shap/capacity/material/share} & \TBDcell{shap/voltage/material/share} & \TBDcell{shap/conductivity/material/share} & \TBDcell{shap/coulombic_efficiency/material/share} & \TBDcell{shap/energy/material/share} \\
        Material name  & \TBDcell{shap/capacity/material_text/share} & \TBDcell{shap/voltage/material_text/share} & \TBDcell{shap/conductivity/material_text/share} & \TBDcell{shap/coulombic_efficiency/material_text/share} & \TBDcell{shap/energy/material_text/share} \\
        Measurement    & \TBDcell{shap/capacity/measurement/share} & \TBDcell{shap/voltage/measurement/share} & \TBDcell{shap/conductivity/measurement/share} & \TBDcell{shap/coulombic_efficiency/measurement/share} & \TBDcell{shap/energy/measurement/share} \\
        \midrule
        Noise ($t_{2,\mathrm{noise}}$) & \TBDcell{shap/capacity/noise/share} & \TBDcell{shap/voltage/noise/share} & \TBDcell{shap/conductivity/noise/share} & \TBDcell{shap/coulombic_efficiency/noise/share} & \TBDcell{shap/energy/noise/share} \\
        \bottomrule
    \end{tabular}
\end{table}

\begin{figure}[htbp]
    \centering
    \FIGTBD{figures/fig_shap_groups_by_target.png}{(a) stacked group shares per target; (b) top-15 features for capacity by mean |SHAP| coloured by group.}
    \caption{\textbf{Grouped SHAP attribution.} (a) Share of total mean $|$SHAP$|$ by feature group for each target ($t_2$, DOI-grouped test records). (b) The 15 highest-ranked predictors for capacity; the suffix \texttt{\_na} denotes a missingness indicator. Values are means over five splits.}
    \label{fig:shap}
\end{figure}

\subsection{Comparison with previous studies on the Huang--Cole database}
\label{sec:previous}

Table~\ref{tab:previous} summarises the studies we found that train machine-learning models on records of the Huang--Cole database. We identified them by screening all 167 works that cite the database in the Semantic Scholar index (November 2026); most citing works are reviews, text-mining tools or databases for other properties. The cell-level and curated studies discussed in the Introduction \cite{severson2019cyclelife,celik2022cyclelife,kauwe2019battery} use different data and are not included here.

Zhuang and Barnard used the database in a series of studies on structure-free composition encodings. They cleaned the formulas of the database into 265,481 instances and clustered them in an unsupervised analysis \cite{zhuang2023mendeleev}, then trained logistic-regression, decision-tree and support-vector classifiers to recover those cluster labels from the same encodings, using a random stratified 85/15 split and 7-fold stratified cross-validation \cite{zhuang2024classification}. Because the labels are clusters of the encoding itself, the near-perfect accuracies reported (0.99--1.00) measure the separability of the encoded compositions rather than the prediction of a measured property; the authors note that the clusters bear no relation to anode or cathode use. A further study classified 471 literature materials by application \cite{zhuang2025applications}. Their regression study \cite{zhuang2024property} is the closest to ours. On 10,129 computed compounds from the Materials Project, formula-only encodings predicted gravimetric capacity, energy and voltage well ($R^2$ up to 0.99) but volumetric properties poorly. They then applied their best combination, the Mendeleev+ encoding (element fractions ordered by Mendeleev number plus statistics of tabulated elemental properties, min--max scaled) with XGBoost tuned over 1,000 hyperparameter combinations by 10-fold cross-validation, to 254,749 Huang--Cole records cleaned with formula-correction rules, using a random 85/15 split. Test $R^2$ was 0.473 for capacity, 0.489 for voltage, 0.495 for energy, 0.142 for Coulombic efficiency and 0.056 for conductivity (modelled in S~cm$^{-1}$ rather than log units). A second variant kept one record per formula with the median value; $R^2$ was nearly unchanged (capacity 0.473), and the authors concluded that removing redundant materials had limited impact. Because each formula then occurs once, a random split of that variant is a composition-grouped split, so their result can be compared directly with our composition-grouped protocol.

Jablonka \cite{jablonka2026clever} used the database as one of five case studies of ``bibliographic confounding''. Capacity records of the merged release without any warning flag (33,701 entries after requiring Crossref author metadata) were featurised with 273 matminer composition descriptors (Magpie element-property statistics, stoichiometry and Meredig descriptors) \cite{ward2018matminer,ward2016magpie,meredig2014combinatorial} and modelled with LightGBM at default settings \cite{ke2017lightgbm}, evaluated by random 10-fold cross-validation. The same descriptors predicted the authors, journal and publication year of a record better than chance, and a second model given only the \emph{predicted} metadata was used to test whether capacity could be predicted through those bibliographic fingerprints; for capacity it was not better than the mean. The study concludes that group and time splits are needed as routine falsification tests, which is the protocol we apply here.

\begin{table*}[t]
    \centering\small
    \caption{\textbf{Previous studies that model records of the Huang--Cole database, and this work.} Values are as reported by each study; tasks, record selection and evaluation differ, so rows are not directly comparable.}
    \label{tab:previous}
    \begin{tabular}{@{}p{2.6cm} p{2.6cm} p{3.0cm} p{2.4cm} p{2.6cm} p{2.2cm}@{}}
        \toprule
        Study & Task & Records & Features / model & Evaluation & Reported \\
        \midrule
        Zhuang \& Barnard \cite{zhuang2024classification} & cluster label (2 or 5 classes) & 265,481 cleaned formula instances & Mendeleev encodings; LR, DT, SVM & random stratified 85/15; 7-fold CV & accuracy 0.99--1.00 \\
        Zhuang \& Barnard \cite{zhuang2024property} & five properties (regression) & 254,749 cleaned records; one-per-formula variant & Mendeleev+ encoding; XGBoost (1,000-combination search) & random 85/15; 10-fold CV & $R^2$ capacity 0.473, voltage 0.489, energy 0.495, CE 0.142, cond.\ 0.056 \\
        Jablonka \cite{jablonka2026clever} & capacity; author, journal, year & 33,701 capacity records (merged release, no warning flag) & 273 matminer descriptors; LightGBM (default) & random 10-fold CV & metadata predictable; metadata-only model $\approx$ mean \\
        \midrule
        This work & five properties & \TBDcell{pop/total/records} records, \TBDcell{pop/total/publications} publications & composition, elemental, name and measurement features; CatBoost--XGBoost & random / DOI-grouped / temporal, 5 repeats & capacity $R^2$ \TBDcell{m/capacity/random/t2/r2} / \TBDcell{m/capacity/doi/t2/r2} / \TBDcell{m/capacity/temporal/t2/r2} \\
        \bottomrule
    \end{tabular}
\end{table*}

To compare on a common basis, we obtained the public implementation of Jablonka \cite{jablonka2026clever} and adapted it to our records and partitions. Their composition parsing (summing the components of \texttt{Extracted\_name} with integer-truncated stoichiometries), featurisation, preprocessing and default LightGBM model were kept unchanged; the model was trained on the training and validation records of split~1 and evaluated on its test records. As in their code, records for which any descriptor cannot be computed (for example formulas that the integer truncation reduces to no element) are removed from training and test; this keeps \TBD{prev/coverage/min}--\TBD{prev/coverage/max} of the test records, depending on the target, and both models are scored on exactly this subset. We adapted the Zhuang--Barnard workflow in the same way. Their regression notebook is part of the supporting information of their article, so we re-implemented it from the Methods section: the Mendeleev+ encoding was built from their periodic table file (atomic mass, group, period, molar volume, atomic number, block and valence; density, atomic volume and Mendeleev number are not in that file), and XGBoost hyperparameters were selected by randomised search with cross-validation on the training records (25 configurations and 3-fold cross-validation on at most 20,000 records, instead of their 1,000 configurations with 10-fold cross-validation). Table~\ref{tab:adapted} reports both adapted pipelines on split~1 for every target under record-level and DOI-grouped splitting. These comparisons evaluate the adapted workflows on our task; they do not reproduce the published scores.

Because Zhuang and Barnard released their cleaned records with DOIs, we could also rerun their workflow on their own data under both split types (Table~\ref{tab:zbB}). On their records and their random 85/15 split the re-implementation gave capacity $R^2$ of \TBD{zbB/capacity/random/r2} (published: 0.473), and a DOI-grouped 85/15 split gave \TBD{zbB/capacity/doi/r2}. The re-implementation reproduced their published random-split results for capacity and voltage to within 0.01 in $R^2$ (energy \TBD{zbB/energy/random/r2} against 0.495), and gave lower values than published for Coulombic efficiency. On the same records, DOI grouping lowered $R^2$ by about 0.06 for capacity and 0.12 for energy and left voltage unchanged, and conductivity and Coulombic efficiency stayed near zero under both splits. This is the same moderate, target-dependent effect we found with our records, features and ensemble, now observed with a different cleaning procedure, encoding, learner and split ratio. The published scores are therefore not strongly inflated by shared publications; the larger gap between published and achievable performance arises from repeated values and new materials (Section~\ref{sec:channels}), which neither split type separates. The energy test set contains only about 1,250 records, so single-split differences of this size for energy should be read with caution.

\begin{table}[htbp]
    \centering\small
    \caption{\textbf{Zhuang--Barnard workflow on their own cleaned records} (re-implementation; random 85/15 split as published, and a DOI-grouped 85/15 split with the same seed). Conductivity is in S~cm$^{-1}$ as in the original study. Published random-split $R^2$ in the first column.}
    \label{tab:zbB}
    \begin{tabular}{@{}lccc@{}}
        \toprule
        Target & Published $R^2$ & Random $R^2$ (re-run) & DOI-grouped $R^2$ \\
        \midrule
        Capacity       & 0.473 & \TBDcell{zbB/capacity/random/r2} & \TBDcell{zbB/capacity/doi/r2} \\
        Voltage        & 0.489 & \TBDcell{zbB/voltage/random/r2} & \TBDcell{zbB/voltage/doi/r2} \\
        Conductivity   & 0.056 & \TBDcell{zbB/conductivity/random/r2} & \TBDcell{zbB/conductivity/doi/r2} \\
        Coulombic eff. & 0.142 & \TBDcell{zbB/coulombic_efficiency/random/r2} & \TBDcell{zbB/coulombic_efficiency/doi/r2} \\
        Energy         & 0.495 & \TBDcell{zbB/energy/random/r2} & \TBDcell{zbB/energy/doi/r2} \\
        \bottomrule
    \end{tabular}
\end{table}

\begin{table*}[t]
    \centering\small
    \caption{\textbf{Adapted pipelines of Jablonka \cite{jablonka2026clever} and Zhuang and Barnard \cite{zhuang2024property} and the CatBoost--XGBoost ensemble on common partitions} (split~1; each adapted pipeline is paired with the ensemble scored on exactly the test records that pipeline can featurise; single-partition values without standard deviations). Cells give $R^2$ / MAE; MAE units as in Table~\ref{tab:split}.}
    \label{tab:adapted}
    \begin{tabular}{@{}llccccc@{}}
        \toprule
        Split & Model & Capacity & Voltage & Conductivity & Coulombic eff. & Energy \\
        \midrule
        Random & Jablonka, adapted & \TBDcell{prev/capacity/random/jab/r2} / \TBDcell{prev/capacity/random/jab/mae} & \TBDcell{prev/voltage/random/jab/r2} / \TBDcell{prev/voltage/random/jab/mae} & \TBDcell{prev/conductivity/random/jab/r2} / \TBDcell{prev/conductivity/random/jab/mae} & \TBDcell{prev/coulombic_efficiency/random/jab/r2} / \TBDcell{prev/coulombic_efficiency/random/jab/mae} & \TBDcell{prev/energy/random/jab/r2} / \TBDcell{prev/energy/random/jab/mae} \\
               & \quad ensemble, same records & \TBDcell{prev/capacity/random/ens/r2} / \TBDcell{prev/capacity/random/ens/mae} & \TBDcell{prev/voltage/random/ens/r2} / \TBDcell{prev/voltage/random/ens/mae} & \TBDcell{prev/conductivity/random/ens/r2} / \TBDcell{prev/conductivity/random/ens/mae} & \TBDcell{prev/coulombic_efficiency/random/ens/r2} / \TBDcell{prev/coulombic_efficiency/random/ens/mae} & \TBDcell{prev/energy/random/ens/r2} / \TBDcell{prev/energy/random/ens/mae} \\
               & Zhuang--Barnard, adapted & \TBDcell{prev/capacity/random/zb/r2} / \TBDcell{prev/capacity/random/zb/mae} & \TBDcell{prev/voltage/random/zb/r2} / \TBDcell{prev/voltage/random/zb/mae} & \TBDcell{prev/conductivity/random/zb/r2} / \TBDcell{prev/conductivity/random/zb/mae} & \TBDcell{prev/coulombic_efficiency/random/zb/r2} / \TBDcell{prev/coulombic_efficiency/random/zb/mae} & \TBDcell{prev/energy/random/zb/r2} / \TBDcell{prev/energy/random/zb/mae} \\
               & \quad ensemble, same records & \TBDcell{prev/capacity/random/ens_zb/r2} / \TBDcell{prev/capacity/random/ens_zb/mae} & \TBDcell{prev/voltage/random/ens_zb/r2} / \TBDcell{prev/voltage/random/ens_zb/mae} & \TBDcell{prev/conductivity/random/ens_zb/r2} / \TBDcell{prev/conductivity/random/ens_zb/mae} & \TBDcell{prev/coulombic_efficiency/random/ens_zb/r2} / \TBDcell{prev/coulombic_efficiency/random/ens_zb/mae} & \TBDcell{prev/energy/random/ens_zb/r2} / \TBDcell{prev/energy/random/ens_zb/mae} \\
        \midrule
        DOI    & Jablonka, adapted & \TBDcell{prev/capacity/doi/jab/r2} / \TBDcell{prev/capacity/doi/jab/mae} & \TBDcell{prev/voltage/doi/jab/r2} / \TBDcell{prev/voltage/doi/jab/mae} & \TBDcell{prev/conductivity/doi/jab/r2} / \TBDcell{prev/conductivity/doi/jab/mae} & \TBDcell{prev/coulombic_efficiency/doi/jab/r2} / \TBDcell{prev/coulombic_efficiency/doi/jab/mae} & \TBDcell{prev/energy/doi/jab/r2} / \TBDcell{prev/energy/doi/jab/mae} \\
               & \quad ensemble, same records & \TBDcell{prev/capacity/doi/ens/r2} / \TBDcell{prev/capacity/doi/ens/mae} & \TBDcell{prev/voltage/doi/ens/r2} / \TBDcell{prev/voltage/doi/ens/mae} & \TBDcell{prev/conductivity/doi/ens/r2} / \TBDcell{prev/conductivity/doi/ens/mae} & \TBDcell{prev/coulombic_efficiency/doi/ens/r2} / \TBDcell{prev/coulombic_efficiency/doi/ens/mae} & \TBDcell{prev/energy/doi/ens/r2} / \TBDcell{prev/energy/doi/ens/mae} \\
               & Zhuang--Barnard, adapted & \TBDcell{prev/capacity/doi/zb/r2} / \TBDcell{prev/capacity/doi/zb/mae} & \TBDcell{prev/voltage/doi/zb/r2} / \TBDcell{prev/voltage/doi/zb/mae} & \TBDcell{prev/conductivity/doi/zb/r2} / \TBDcell{prev/conductivity/doi/zb/mae} & \TBDcell{prev/coulombic_efficiency/doi/zb/r2} / \TBDcell{prev/coulombic_efficiency/doi/zb/mae} & \TBDcell{prev/energy/doi/zb/r2} / \TBDcell{prev/energy/doi/zb/mae} \\
               & \quad ensemble, same records & \TBDcell{prev/capacity/doi/ens_zb/r2} / \TBDcell{prev/capacity/doi/ens_zb/mae} & \TBDcell{prev/voltage/doi/ens_zb/r2} / \TBDcell{prev/voltage/doi/ens_zb/mae} & \TBDcell{prev/conductivity/doi/ens_zb/r2} / \TBDcell{prev/conductivity/doi/ens_zb/mae} & \TBDcell{prev/coulombic_efficiency/doi/ens_zb/r2} / \TBDcell{prev/coulombic_efficiency/doi/ens_zb/mae} & \TBDcell{prev/energy/doi/ens_zb/r2} / \TBDcell{prev/energy/doi/ens_zb/mae} \\
        \bottomrule
    \end{tabular}
\end{table*}

\textbf{Bibliographic fingerprints under DOI grouping.} We also repeated the metadata test of Jablonka with the metadata that the database itself provides (journal, as the ten most frequent journals, and publication year; author lists are not part of the release). From the matminer descriptors, publication year was predicted with an MAE of \TBD{ch/capacity/random/year_mae} years under record-level splitting and \TBD{ch/capacity/doi/year_mae} years under DOI grouping, against \TBD{ch/capacity/doi/year_mae_baseline} years for the training mean. A model given only the predicted journal and year reached a capacity $R^2$ of \TBD{ch/capacity/random/indirect_r2} (random) and \TBD{ch/capacity/doi/indirect_r2} (DOI-grouped). The same was true for every target (Supplementary Table S1). Composition predicted the journal noticeably better than chance only under record-level splitting (micro-F1 \TBD{ch/capacity/random/journal_f1} versus \TBD{ch/capacity/doi/journal_f1} for capacity), that is, by recognising publications that were also in training. Jablonka found descriptor-based ``bibliographic fingerprints'' in several databases; in this one they are weak and are mostly a symptom of shared publications, which is consistent with the small record-level advantage reported above.

\textbf{Implications for evaluating models on text-mined data.} Publication-grouped evaluation is usually motivated by four assumptions: records from one publication are correlated; record-level splits let a model exploit that correlation; features that are constant within a publication are especially prone to it; and a DOI-grouped split removes it. In the battery database the first holds strongly (ICC(1) 0.30--0.84), but the second holds only weakly and the fourth only partly.

\textit{Features decide whether clustering becomes leakage.} Device and process databases carry many variables whose joint values are specific to a laboratory; the battery records carry almost none, and their composition descriptors are shared by hundreds of publications. Clustered targets alone did not produce a large record-level advantage. A model can only exploit publication identity that its inputs encode, which is why reporting the between-publication variance is not enough and the random-versus-grouped comparison has to be run for each feature set.

\textit{DOI grouping does not remove cross-publication repetition.} In a curated device database each record is a new device. In a text-mined materials database a large share of values are restatements of values reported elsewhere, and grouping by DOI cannot separate them. Composition-grouped and temporal splits, and the repeated/non-repeated breakdown, are needed to estimate how a model performs on values that have no precedent in the literature, which is the case that matters for discovery.

\textit{Long format and missing conditions.} A record is a single number, not a device, so the leakage argument has to be made per target, and the absence of electrode formulation, electrolyte, loading and voltage window caps achievable accuracy, most visibly for Coulombic efficiency. These limits are properties of the database, not of the model: the adapted pipeline of Jablonka, with different descriptors and learner, shows the same pattern across targets and protocols (Table~\ref{tab:adapted}).

\textit{Label noise is structured.} With an extraction precision of about 80\%, errors are themselves clustered by publication (they depend on how a paper is written) and repeated values can propagate an error across publications. Part of what any split rewards is the reproduction of the database's own extraction behaviour.

\subsection{Limitations}
\label{sec:limitations}

DOI-grouped evaluation tests performance on publications excluded from training; it does not guarantee independence of laboratories or research groups, and author information is not part of the release. Our repeated-value analysis matches on the normalised compound name and the exact value, so it misses restatements with rounding or a different name and undercounts recurrence. Composition grouping uses rounded atomic fractions of the parsed formula and therefore treats composites with the same overall composition as one material. The targets mix quantities that a materials scientist would keep apart (theoretical and measured capacity, average and cut-off voltage, ionic and electronic conductivity), and the specifier flags only partly separate them. Material descriptors are computed from the parsed formula of the active material and ignore binders, conductive additives, electrolytes and counter electrodes \cite{kauwe2019battery}. The database ends with papers published in \TBD{pop/date_max_year}, so the temporal split cannot test generalisation to recent chemistries. The adapted pipelines keep the published featurisation and learner but are evaluated on our records and partitions, so they test the workflows on a common task and do not reproduce the published scores. Finally, the extraction errors discussed in Section~\ref{sec:quality} are not removed by any of our filters.

% ---------------- CONCLUSION ----------------
\section{Conclusion}
\label{sec:conclusion}

We trained CatBoost--XGBoost ensembles on the Huang--Cole auto-generated battery database for five properties and evaluated identical models under record-level, DOI-grouped, temporal and composition-grouped splits. Records are strongly clustered by publication, yet DOI grouping lowered $R^2$ by only 0.02--0.08 (capacity \TBD{m/capacity/random/t2/r2} to \TBD{m/capacity/doi/t2/r2}), because composition descriptors carry little recoverable publication identity. A larger effect comes from values that recur across publications: \TBD{rep/doi/share-range} of DOI-grouped test records have such a precedent and are predicted much better than the rest, and on values without one capacity $R^2$ falls to \TBD{rep/capacity/doi/r2_new} and conductivity and Coulombic efficiency are not predicted at all. Temporal and composition-grouped splits lower performance further. Title keywords, which describe the cell chemistry, improve predictions under both split types, the raw compound name changes DOI-grouped $R^2$ by at most 0.01, Gaussian-noise columns receive \TBD{shap/noise/share} of the attribution and do not improve any target, and removing within-paper series narrows the record-level advantage further.

Three recommendations follow for models trained on text-mined materials data. First, report publication-grouped results, but do not treat them as a test of generalisation to new materials: add composition-grouped or temporal splits and report performance separately for values with and without a precedent in the training literature. Second, check whether the features can identify publications, for example with the bibliographic-fingerprint test, before interpreting a small or large random-versus-grouped gap. Third, audit the database's own flags before use, because in this database one flag encodes publication identity and another is a function of the target. Future work should add the experimental context that the database lacks, separate theoretical from measured values at extraction time, and test whether newer extraction pipelines \cite{huang2022batterybert,dagdelen2024structured} change these conclusions.

\section*{Author contributions}
S.~Patel: conceptualisation, methodology, software, formal analysis, writing -- original draft. S.~Misra: supervision, conceptualisation, writing -- review and editing.

\section*{Conflicts of interest}
There are no conflicts to declare.

\section*{Data availability}
The Huang--Cole battery database is available from figshare \cite{huang2020figshare}. The code for data preparation, model training, evaluation and the adapted comparison pipeline, the frozen element list, the split fingerprints and all result tables are available at \url{https://github.com/sunraj-patel/battery-prediction}; it will be archived at Zenodo with a DOI on acceptance. The adapted pipelines reuse code from \url{https://github.com/lamalab-org/clever-materials-hans} (MIT licence) and data files from \url{https://github.com/ZoraZhuang/battery_datasets} (GPL-3.0).

\section*{Declaration of generative AI use}
During this work the authors used Claude (Anthropic) to write and adapt the analysis code, to search and verify the literature, and to draft and revise the manuscript text. All code was run and all results were checked by the authors, who take full responsibility for the content.

\bibliographystyle{elsarticle-num}
\bibliography{references_battery}

\end{document}
